\section{Main Structural Theorem}

We now collect the native register, quotient, matching, cycle, and
extension results into one statement.

\begin{theorem}[Native Klein register theorem]
Let
\[
G_{60}\cong \operatorname{AT4val}[60,6]
\]
be the exact connected quartic graph considered throughout this paper.
Then the following statements hold.

\begin{enumerate}
    \item The graph contains a free normal Klein four-subgroup
    \[
    V_4^{\mathrm{nat}}
    =
    \{1,a,b,ab\}
    \trianglelefteq
    \operatorname{Aut}(G_{60}).
    \]

    \item The native Klein action has fifteen orbits of size four:
    \[
    V(G_{60})
    =
    \bigsqcup_{i=0}^{14} C_i,
    \qquad
    |C_i|=4.
    \]

    \item The quotient by the native Klein group is the Petersen line
    graph:
    \[
    G_{60}/V_4^{\mathrm{nat}}
    \cong
    L(P).
    \]

    \item The native state space decomposes into four orthogonal
    fifteen-dimensional character sectors:
    \[
    \mathbb R^{60}
    =
    \mathcal H_{\chi_1}
    \oplus
    \mathcal H_{\chi_a}
    \oplus
    \mathcal H_{\chi_b}
    \oplus
    \mathcal H_{\chi_{ab}}.
    \]

    \item The full automorphism group fixes \(a\) and exchanges \(b\)
    with \(ab\) under exactly the odd half of the canonical
    five-point action. Consequently,
    \[
    \chi_a
    \mid
    \{\chi_b,\chi_{ab}\},
    \]
    and the \(\chi_b\)- and \(\chi_{ab}\)-operators are
    symmetry-protected isospectral sectors.

    \item The fifteen chambers admit an invariant partition
    \[
    \mathcal M
    =
    \{T_0,T_1,T_2,T_3,T_4\}
    \]
    into five blocks of size three.

    \item Under the quotient identification with \(L(P)\), each
    block \(T_i\) is an independent three-set and corresponds to a
    three-edge matching \(M_i\) in the Petersen graph. The five
    matchings partition the Petersen edge set:
    \[
    E(P)
    =
    M_0\sqcup M_1\sqcup M_2\sqcup M_3\sqcup M_4.
    \]

    \item The induced action on the five matching blocks is the full
    symmetric group:
    \[
    \rho:
    \operatorname{Aut}(G_{60})
    \twoheadrightarrow
    S_5,
    \]
    with
    \[
    \ker\rho
    =
    V_4^{\mathrm{nat}}.
    \]

    \item The resulting exact sequence splits:
    \[
    1
    \longrightarrow
    V_4^{\mathrm{nat}}
    \longrightarrow
    \operatorname{Aut}(G_{60})
    \longrightarrow
    S_5
    \longrightarrow
    1.
    \]

    \item The conjugation action factors through the sign character:
    \[
    \operatorname{Aut}(G_{60})
    \cong
    V_4\rtimes_{\operatorname{sgn}}S_5,
    \]
    where every element of \(S_5\) fixes \(a\), every even element
    fixes \(b\) and \(ab\) separately, and every odd element
    exchanges them.

    \item The semidirect-product description is equivalent to the
    previously established fiber-product classification:
    \[
    V_4\rtimes_{\operatorname{sgn}}S_5
    \cong
    S_5\times_{C_2}D_8.
    \]

    \item The shortest quotient cycles with nontrivial native
    holonomy are the twelve quotient pentagons. Six have holonomy
    \(b\), and six have holonomy \(ab\).

    \item Each quotient pentagon lifts to two disjoint native
    ten-cycles. These twenty-four registered ten-cycles are exactly
    the simple ten-cycles whose quotient projection has minimal
    period five.

    \item The twenty-four registered cycles form one automorphism
    orbit. Their individual stabilizers and answering-pair
    stabilizers are
    \[
    D_{10}
    \quad\text{and}\quad
    D_{10}\times C_2,
    \]
    respectively.

    \item Registered ten-cycles admit the exact fourth-order
    recognition laws
    \[
    x_C^{\mathsf T}A^4x_C=880
    \]
    and
    \[
    x_C^{\mathsf T}L^4x_C=2160.
    \]

    \item The adjacent third-moment frontier contains \(120\)
    nonregistered cycles. Each carries two external relay paths
    forming
    \[
    P_3\sqcup P_3
    \]
    and exactly sixteen additional external four-step returns.

    \item The \(120\) frontier cycles are stabilized by fifteen
    distinct foreign Klein four-groups. Each foreign subgroup
    stabilizes eight cycles, intersects the native Klein group only
    in the identity, and is not conjugate to it.

    \item The chamber-incidence graph of the fifteen foreign Klein
    groups is
    \[
    5K_3.
    \]
    Its five components are exactly the five matching blocks
    \(T_0,\ldots,T_4\).
\end{enumerate}
\end{theorem}

\subsection{Proof from the native register}

The native involutions \(a\), \(b\), and \(ab\) are reconstructed as
fixed-point-free graph automorphisms and generate a free Klein
four-group. Direct conjugation by all \(480\) native automorphisms
proves normality and yields the intrinsic conjugation grammar
\[
a\mid\{b,ab\}.
\]

The free action gives fifteen four-vertex chambers. The exact orbit
quotient is quartic on fifteen vertices and is identified with
\(L(P)\). The regular representation of \(V_4\) on each chamber
produces the four equal character sectors. Gauge switching changes
the signed quotient matrices by diagonal conjugation, so their spectra
are intrinsic.

The foreign-Klein stabilizer census then supplies a second relation
on the same fifteen chambers. Its connected components are five
three-sets. These sets are independent in the quotient carrier and
therefore become three-edge matchings under the Petersen-edge bridge.

\subsection{Proof from the five-block action}

Every one of the \(480\) automorphisms preserves the five matching
blocks. The induced action contains \(120\) distinct block
permutations and is therefore the full group \(S_5\).

The kernel has order four and agrees exactly with
\(V_4^{\mathrm{nat}}\). This proves
\[
\operatorname{Aut}(G_{60})/V_4^{\mathrm{nat}}
\cong
S_5.
\]

A bounded lift search for a transposition and a five-cycle finds
subgroups of order \(120\) intersecting the native kernel trivially.
Hence the extension splits. The complete conjugation census then
shows that the quotient action on the kernel is exactly the sign
action.

\subsection{Compatibility with the earlier fiber-product theorem}

The earlier companion paper proves
\[
\operatorname{Aut}(G_{60})
\cong
S_5\times_{C_2}D_8
\]
through internal group invariants and a complete explicit
multiplication certificate.

The present theorem does not replace that classification. It gives a
native realization of it.

The kernel of the projection
\[
S_5\times_{C_2}D_8
\longrightarrow
S_5
\]
is the Klein four-group
\[
\ker\chi_{V_4}\leq D_8.
\]

An element of \(D_8\) outside this kernel acts on it by fixing one
nonidentity involution and exchanging the other two. The fiber-product
condition couples this action to the sign of the \(S_5\)-coordinate.

Thus
\[
S_5\times_{C_2}D_8
\]
and
\[
V_4\rtimes_{\operatorname{sgn}}S_5
\]
are two descriptions of the same group architecture.

The former is adapted to abstract classification. The latter is
adapted to the native graph:

\[
\boxed{
5\text{ matching blocks}
\times
3\text{ chambers per block}
\times
4\text{ native Klein states per chamber}.
}
\]

\subsection{Scope of the theorem}

The theorem establishes a finite graph-theoretic structure. It
identifies:

\begin{enumerate}
    \item the native hidden kernel;
    \item the visible five-point quotient;
    \item the parity coupling between them;
    \item the quotient character operators;
    \item the registered-cycle class;
    \item the foreign stabilizer geometry.
\end{enumerate}

It does not identify these quantities with physical energy, mass,
charge, spin, time, spacetime dimensions, or any physical interaction.

The earned conclusion is structural:

\begin{quote}
The native sixty-vertex graph is an equivariant
\(5\times3\times4\) register whose visible \(S_5\)-motion is coupled
to a hidden Klein state through permutation parity.
\end{quote}
